% Eksplorator fraktali - zarys projektu do akceptacji (wersja 3, zgodna z opisem projektu)
% Kompilacja: pdflatex zarys_projektu_v2.tex (dwa razy)

\documentclass[11pt,a4paper]{article}

\usepackage[utf8]{inputenc}
\usepackage[T1]{fontenc}
\usepackage[polish]{babel}
\usepackage{lmodern}
\usepackage[margin=2.2cm]{geometry}
\usepackage{amsmath,amssymb}
\usepackage{enumitem}
\usepackage{booktabs}
\usepackage{tabularx}
\usepackage{array}
\usepackage{xcolor}
\usepackage{graphicx}
\usepackage{fancyvrb}
\usepackage{tikz}
\usetikzlibrary{positioning}
\usepackage[hidelinks]{hyperref}

\setlist{itemsep=2pt,topsep=4pt}
\setlength{\parskip}{4pt}
\setlength{\parindent}{0pt}
\setlength{\emergencystretch}{3em}

\newcommand{\kod}[1]{\texttt{#1}}
\fvset{frame=single,fontsize=\small,framesep=6pt}

\begin{document}

\begin{titlepage}
    \centering
    \IfFileExists{AGH.pdf}
      {\includegraphics[width=0.3\textwidth]{AGH.pdf}}
      {\fbox{\parbox[c][4cm][c]{0.3\textwidth}{\centering logo AGH}}}\\[1cm]
    \textsc{\LARGE Akademia Górniczo-Hutnicza\\
    im. Stanisława Staszica w Krakowie}\\[2cm]
    \textsc{\Large Wydział Fizyki i Informatyki Stosowanej}\\[1cm]
    \textsc{\large Interfejsy użytkownika i biblioteki graficzne}\\[1cm]
    {\huge \bfseries Eksplorator fraktali}\\[0.4cm]
    {\Large Zarys projektu do akceptacji}\\[2cm]
    {\Large \textbf{Autorzy}}\\[0.5cm]
    {\large Jakub Piela (kapitan)\\[2pt] Marcel Gładysz\\[2pt] Patryk Pilch}\\
    \vfill
    \begin{flushright}
    \textbf{Prowadzący:} dr inż. Janusz Malinowski\\
    \textbf{Temat:} własny\\
    \textbf{Rok akademicki:} 2026/2027\\
    \textbf{Data złożenia:} 02.11.2026\\
    \end{flushright}
\end{titlepage}

% ---------------------------------------------------------------
\section{Opis projektu}

Eksplorator fraktali to program w C++ do generowania, interaktywnego oglądania
i animowania fraktali z trzech rodzin:
\begin{itemize}
  \item \textbf{fraktale zespolone} -- zbiór Mandelbrota i zbiory Julii,
  \item \textbf{fraktale IFS} -- atraktory układów przekształceń afinicznych,
        np.\ trójkąt Sierpińskiego i paproć Barnsleya,
  \item \textbf{L-systemy} -- gramatyki przepisywania napisów rysowane grafiką
        żółwia, np.\ krzywa Kocha i rośliny.
\end{itemize}

Użytkownik przybliża i przesuwa widok myszą, zmienia parametry i palety
kolorów w panelu sterowania, a bieżący obraz zapisuje do pliku. Program
tworzy też \textbf{animację płynnego przejścia (,,morfingu'') między
fraktalami IFS}: przeliczając współczynniki przekształceń z jednego fraktala na
drugi, generuje serię bitmap zapisywanych na dysku.

Każdy piksel obrazu liczy nasz kod na procesorze. Biblioteki służą wyłącznie
do otwarcia okna, obsługi zdarzeń, wyświetlenia gotowego bufora pikseli
i narysowania panelu sterowania.

Animacja przejścia oraz -- w wersji rozszerzonej -- plik scenariusza
i fraktale trójwymiarowe nawiązują do niewybranego tematu nr~25 (,,Generowanie
i animacja fraktali''). Nasz projekt obejmuje jego zakres dotyczący fraktali
IFS i rozszerza go o fraktale zespolone oraz L-systemy.

% ---------------------------------------------------------------
\section{Założenia wstępne}

\begin{itemize}
  \item \textbf{Język i środowisko:} C++17, Visual Studio 2026, Windows~11,
        kompilacja x64. Projekt budujemy przez CMake, więc kompiluje się też
        na Linuksie. SFML jest linkowany statycznie, więc program działa bez
        dodatkowych plików DLL.
  \item \textbf{Biblioteki z wykładu:} SFML~3.1 (okno, zdarzenia, wyświetlenie
        bufora pikseli), Dear ImGui z nakładką ImGui-SFML (panel sterowania)
        oraz ImGuiFileDialog (okna wyboru plików).
  \item \textbf{Bez OpenGL i DirectX w naszym kodzie:} obliczenia, kolorowanie,
        rysowanie linii i rzutowanie fraktali 3D wykonuje nasz kod na CPU.
  \item \textbf{Działanie bez plików wejściowych:} wszystkie parametry można
        ustawić w panelu. Pliki definicji są opcjonalne i służą do zapisu oraz
        wczytywania gotowych ustawień.
  \item \textbf{Jeden wątek w wersji podstawowej:} obliczenia wykonuje główny
        wątek programu. Animację generujemy po jednej klatce na obieg pętli
        głównej, więc między klatkami okno się odświeża i pokazuje podgląd.
        Obliczenia wielowątkowe i generowanie animacji w tle są poszerzeniem.
  \item \textbf{Okno} jest skalowalne. Po zmianie rozmiaru obraz jest
        przeliczany w nowej rozdzielczości, a nie rozciągany.
  \item \textbf{Oddanie projektu:} archiwum \kod{numer\_projektu.7z}
        (kompresja Ultra) z katalogami \kod{DOC}, \kod{SRC} i \kod{BIN};
        dokumentacja w PDF zgodna ze wzorcem.
\end{itemize}

% ---------------------------------------------------------------
\section{Wymagania}

\subsection{Wymagania podstawowe}

\begin{enumerate}
  \item \textbf{Interfejs graficzny:} okno z obszarem fraktala, panel
        sterowania z polami i suwakami do zmiany parametrów, menu oraz okna
        wyboru plików. Po zmianie rozmiaru okna obraz jest przeliczany w nowej
        rozdzielczości, a nie rozciągany.
  \item \textbf{Zbiór Mandelbrota:} przybliżanie kółkiem myszy w punkt pod
        kursorem, przesuwanie przeciąganiem, powrót do widoku startowego,
        regulowana maksymalna liczba iteracji.
  \item \textbf{Zbiory Julii:} parametr $c$ wybierany kliknięciem na zbiorze
        Mandelbrota albo wpisywany w panelu.
  \item \textbf{Kolorowanie zbiorów zespolonych:} płynne (bez widocznych pasów),
        co najmniej trzy palety do wyboru.
  \item \textbf{Fraktale IFS:} rysowanie grą w chaos; wbudowane trójkąt
        Sierpińskiego i paproć Barnsleya; edycja współczynników w panelu oraz
        wczytywanie i zapis definicji w pliku tekstowym. Kolor punktu zależy
        od przekształcenia, którym go wyznaczono, a fraktal jest automatycznie
        skalowany tak, aby mieścił się w oknie i nie był za mały.
  \item \textbf{Animacja przejścia IFS:} płynne przejście między dwoma
        fraktalami IFS w zadanej przez użytkownika liczbie klatek, podgląd
        klatek w trakcie generowania, zapis ponumerowanej serii obrazów PNG
        lub BMP.
  \item \textbf{L-systemy:} aksjomat, reguły, kąt i liczba iteracji z panelu
        lub pliku; rysowanie grafiką żółwia ze stosem rozgałęzień; wbudowane
        krzywa Kocha, smok Heighwaya i roślina; automatyczne dopasowanie
        rysunku do okna.
  \item \textbf{Zapis obrazu:} bieżący widok do pliku PNG lub BMP.
\end{enumerate}

\subsection{Wymagania poszerzone}

\begin{itemize}
  \item \textbf{Obliczenia wielowątkowe:} zbiory zespolone liczone na wielu
        wątkach, aby przybliżanie i przesuwanie były płynne; animacja
        generowana w tle z paskiem postępu i możliwością przerwania.
  \item \textbf{Plik scenariusza:} animacja przez więcej niż dwa fraktale
        (kolejne fraktale i liczba klatek każdego przejścia).
  \item \textbf{Fraktale IFS 3D:} przekształcenia przestrzenne (12
        współczynników), własne rzutowanie perspektywiczne i pozycja
        obserwatora ustawiana w panelu. Punkty dalsze są ciemniejsze i nie
        przesłaniają bliższych (własny bufor głębi).
  \item \textbf{Animacja przybliżania} zbioru Mandelbrota lub Julii
        eksportowana jako seria klatek.
  \item \textbf{Edytor palet:} gradient definiowany punktami kontrolnymi, zapis
        i wczytywanie palet.
  \item \textbf{Kolejne fraktale zespolone:} Burning Ship, Multibrot ($z^n$),
        fraktal Newtona dla wielomianu $z^3 - 1$.
  \item \textbf{Zapis w wysokiej rozdzielczości} (większej niż okno)
        z wygładzaniem krawędzi (supersampling).
\end{itemize}

% ---------------------------------------------------------------
\section{Analiza projektu}

\subsection{Specyfikacja danych wejściowych}

\begin{center}
\begin{tabularx}{\textwidth}{@{}l l X@{}}
  \toprule
  Rodzaj & Format & Zawartość \\
  \midrule
  Interakcja & mysz, klawiatura, panel &
    przybliżanie, przesuwanie, wybór $c$, parametry fraktali, paleta \\
  Definicja IFS & plik \kod{.ifs} &
    jedno przekształcenie w wierszu: \kod{a b c d e f p} \\
  L-system & plik \kod{.lsys} &
    aksjomat, reguły w postaci \kod{F=F+F-F}, kąt, liczba iteracji \\
  Scenariusz animacji & plik \kod{.anim} &
    (poszerzenie) rozmiar klatek, liczba punktów, kolejne fraktale, liczba
    klatek każdego przejścia \\
  \bottomrule
\end{tabularx}
\end{center}

Przekształcenie IFS ma postać $x' = ax + by + e$, $y' = cx + dy + f$, a $p$ to
prawdopodobieństwo, z jakim jest losowane w grze w chaos. Jeśli
prawdopodobieństwa w pliku nie sumują się do 1, program je normalizuje.
We wszystkich plikach znak \kod{\#} zaczyna komentarz, a puste wiersze są
pomijane. Przykład -- trójkąt Sierpińskiego:

\begin{Verbatim}
# a b c d e f p
0.5 0 0 0.5 0    0    0.3333
0.5 0 0 0.5 0.5  0    0.3333
0.5 0 0 0.5 0.25 0.5  0.3334
\end{Verbatim}

Plik L-systemu:
\begin{Verbatim}
axiom F
rule F=F+F-F-F+F
angle 90
iterations 4
\end{Verbatim}

Plik scenariusza (poszerzenie) zawiera rozmiar klatek i liczbę punktów na
klatkę, a potem kolejne fraktale w formacie \kod{.ifs} poprzedzone liczbą
przekształceń, przedzielone liczbą klatek przejścia:

\begin{Verbatim}
# szerokosc wysokosc punkty_na_klatke
800 600 2000000
# liczba fraktali
2
# fraktal 1: liczba przeksztalcen, potem wiersze a b c d e f p
3
0.5 0 0 0.5 0    0    0.3333
0.5 0 0 0.5 0.5  0    0.3333
0.5 0 0 0.5 0.25 0.5  0.3334
# liczba klatek przejscia do fraktala 2
100
# fraktal 2
4
 0.00  0.00  0.00  0.16  0  0     0.01
 0.85  0.04 -0.04  0.85  0  1.6   0.85
 0.20 -0.26  0.23  0.22  0  1.6   0.07
-0.15  0.28  0.26  0.24  0  0.44  0.07
\end{Verbatim}

Błędne pliki (zła liczba wartości, niepoprawne liczby, brak sekcji) nie
zamykają programu. Wyświetlamy komunikat z numerem wiersza, w którym wystąpił
błąd.

\subsection{Opis oczekiwanych danych wyjściowych}

\begin{itemize}
  \item \textbf{Obraz:} bieżący widok zapisany jako PNG lub BMP w rozmiarze
        okna albo w zadanej rozdzielczości.
  \item \textbf{Animacja:} seria ponumerowanych plików
        \kod{klatka\_0000.png}, \kod{klatka\_0001.png}, \ldots\ w wybranym
        katalogu. Kolejno oglądane tworzą płynną animację.
  \item \textbf{Definicje:} bieżące ustawienia IFS lub L-systemu zapisane do
        pliku w formatach opisanych powyżej.
\end{itemize}

\subsection{Struktury danych}

\begin{center}
\begin{tabularx}{\textwidth}{@{}l X@{}}
  \toprule
  Struktura & Zawartość \\
  \midrule
  \kod{Color} & kolor RGBA, 4 bajty (układ zgodny z \kod{sf::Texture}
                i \kod{sf::Image}) \\
  \kod{Image} & bufor pikseli \kod{std::vector<Color>} oraz szerokość
                i wysokość; zapis piksela i odcinka (Bresenham) \\
  \kod{LineReader} & czytnik plików tekstowych wspólny dla parserów: pomija
                     komentarze i puste wiersze, pamięta numer wiersza \\
  \kod{ParseError} & wyjątek z numerem wiersza i komunikatem błędu \\
  \kod{ComplexView} & środek widoku $(c_x, c_y)$, skala (jednostki na piksel),
                      maksymalna liczba iteracji \\
  \kod{Palette} & lista punktów kontrolnych gradientu (pozycja $t$, kolor) \\
  \kod{AffineMap} & współczynniki $a, b, c, d, e, f$ i prawdopodobieństwo $p$ \\
  \kod{Ifs} & nazwa i \kod{std::vector<AffineMap>} \\
  \kod{AffineMap3D} & macierz $3 \times 3$, wektor przesunięcia i $p$
                      (poszerzenie) \\
  \kod{Morph} & dwa fraktale IFS, liczba klatek, numer bieżącej klatki,
                katalog i format zapisu \\
  \kod{AnimationScenario} & (poszerzenie) rozmiar klatek, liczba punktów na
                            klatkę, kolejne fraktale, liczby klatek przejść \\
  \kod{LSystem} & aksjomat, reguły \kod{std::map<char, std::string>}, kąt,
                  liczba iteracji \\
  \kod{TurtleState} & pozycja $(x, y)$ i kierunek żółwia; rozgałęzienia
                      na \kod{std::stack<TurtleState>} \\
  \bottomrule
\end{tabularx}
\end{center}

Każdy rodzaj fraktala implementuje wspólną klasę bazową \kod{FractalRenderer}.
Klasy pochodne to \kod{MandelbrotRenderer}, \kod{JuliaRenderer},
\kod{IfsRenderer} i \kod{LSystemRenderer}. Aplikacja zna tylko ten interfejs,
więc zmiana trybu sprowadza się do wyboru innego obiektu:

\begin{itemize}
  \item \kod{render(Image\&)} -- liczy i rysuje obraz; nie korzysta z ImGui
        ani SFML,
  \item \kod{drawControls()} -- rysuje widżety trybu w panelu sterowania,
  \item \kod{tick()} -- wywoływana raz na obieg pętli głównej; służy do pracy
        rozłożonej na wiele klatek, np.\ generowania animacji,
  \item \kod{onZoom}, \kod{onDrag}, \kod{onClick}, \kod{resetView} --
        nawigacja myszą i powrót do widoku startowego,
  \item \kod{loadDefinition}, \kod{saveDefinition} -- wczytywanie i zapis
        plików \kod{.ifs} i \kod{.lsys}; błędy zgłaszane przez
        \kod{ParseError}.
\end{itemize}

Metody zwracają \kod{true}, gdy obraz trzeba przeliczyć. Aplikacja wywołuje
wtedy \kod{render()}.

\subsection{Specyfikacja interfejsu użytkownika}

\begin{center}
\begin{tikzpicture}[font=\small]
  \draw[thick, rounded corners=2pt] (0,0) rectangle (14,7.8);
  \draw (0,7.3) -- (14,7.3);
  \node[anchor=west] at (0.15,7.55) {\textbf{Plik}\quad \textbf{Widok}\quad \textbf{Pomoc}};
  \draw (0,0.45) -- (14,0.45);
  \node[anchor=west] at (0.15,0.22)
    {\footnotesize $x = -0.7436$,\ \ $y = 0.1318$\quad skala: $2.1\cdot10^{-5}$\quad czas obliczeń: 38 ms};
  \fill[gray!12] (0,0.45) rectangle (9.2,7.3);
  \node[align=center] at (4.6,3.9) {\textbf{Obszar fraktala}\\[2pt]
    kółko -- przybliżanie w punkt pod kursorem\\
    przeciąganie -- przesuwanie\\
    prawy przycisk -- wybór $c$ dla zbioru Julii};
  \draw (9.2,0.45) -- (9.2,7.3);
  \node[anchor=north west, align=left, font=\footnotesize] at (9.35,7.15) {
    \small\textbf{Tryb:} Mandelbrot $\blacktriangledown$\\[4pt]
    \small\textbf{Parametry}\\
    iteracje, $c$, współczynniki IFS,\\
    reguły L-systemu\\[4pt]
    \small\textbf{Paleta}\\
    wybór palety, podgląd\\[4pt]
    \small\textbf{Animacja}\\
    fraktal A i B,\\
    liczba klatek, katalog\\
    {[}Generuj{]}\ \ {[}Przerwij{]}\\
    pasek postępu\\[4pt]
    \small\textbf{Zapis}\\
    {[}Zapisz obraz{]}\\
    {[}Zapisz definicję{]}};
\end{tikzpicture}
\end{center}

Okno składa się z menu (wczytywanie i zapis plików, powrót do widoku
startowego, informacje o programie), obszaru fraktala, panelu sterowania
po prawej i paska stanu. Pasek stanu pokazuje współrzędne punktu pod kursorem,
bieżącą skalę i czas obliczenia klatki. Elementy, których nie da się w danej
chwili użyć, są wyszarzone, np.\ zapis definicji w trybie, który jej nie ma.
Pasek postępu animacji pokazuje numer bieżącej klatki, a przycisk
\textit{Przerwij} zatrzymuje generowanie po bieżącej klatce. Panel i okna
dialogowe tworzymy w Dear ImGui.

\subsection{Wyodrębnienie i zdefiniowanie zadań}

\begin{enumerate}
  \item \textbf{Rdzeń aplikacji i interfejs:} okno, pętla zdarzeń, panel ImGui,
        nawigacja myszą, okna dialogowe, zmiana rozmiaru okna, integracja
        modułów.
  \item \textbf{Fraktale zespolone:} Mandelbrot i Julia, płynne kolorowanie,
        palety.
  \item \textbf{IFS i animacja:} parser plików, gra w chaos, kolorowanie,
        dopasowanie do okna, przejście między dwoma fraktalami, eksport
        ponumerowanych klatek.
  \item \textbf{L-systemy:} parser, przepisywanie napisów, grafika żółwia,
        rysowanie odcinków w buforze pikseli.
  \item \textbf{Część wspólna:} klasy \kod{Image}, \kod{LineReader},
        \kod{ParseError} i interfejs \kod{FractalRenderer}, zapis PNG i BMP,
        konfiguracja projektu (CMake, submoduły, formatowanie i analiza kodu),
        testy.
\end{enumerate}

\subsection{Decyzja o wyborze narzędzi programistycznych}

\begin{itemize}
  \item \textbf{Visual Studio 2026 i C++17} -- znamy je z zajęć. Standard
        daje \kod{<random>} do losowania w grze w chaos, \kod{<filesystem>} do
        zapisu klatek oraz \kod{std::thread} i \kod{std::atomic} do poszerzeń
        wielowątkowych.
  \item \textbf{CMake} -- jeden opis budowania dla Visual Studio i Linuksa.
        Biblioteki są dołączone jako submoduły git w ustalonych wersjach, więc
        każdy w zespole buduje ten sam kod. Każdy moduł (\kod{core},
        \kod{complex}, \kod{ifs}, \kod{lsystem}, \kod{app}) to osobny cel,
        co ułatwia równoległą pracę i testy bez okna (CTest).
  \item \textbf{clang-format i clang-tidy} -- wspólny styl kodu
        i analiza statyczna. clang-tidy sprawdza tylko nasz kod, nie
        dołączone biblioteki.
  \item \textbf{SFML 3.1} -- poznana na wykładzie i laboratoriach, prosta
        obsługa okna, zdarzeń i wyświetlania bufora przez \kod{sf::Texture}.
        \kod{sf::Image} zapisuje obrazy PNG i BMP.
  \item \textbf{Dear ImGui, ImGui-SFML, ImGuiFileDialog} -- omawiane na
        wykładzie. Pozwalają szybko zbudować panel z suwakami, polami liczbowymi
        i oknami wyboru plików.
  \item \textbf{Git} (wspólne repozytorium) do pracy zespołowej, \LaTeX\ do
        dokumentacji, 7-Zip do archiwum.
\end{itemize}

% ---------------------------------------------------------------
\section{Algorytmy}

\subsection{Zbiory Mandelbrota i Julii}

Dla każdego punktu płaszczyzny zespolonej powtarzamy przekształcenie
\[
  z_{n+1} = z_n^2 + c .
\]
W zbiorze Mandelbrota $z_0 = 0$, a $c$ jest współrzędną piksela. W zbiorze
Julii $c$ jest stałe dla całego obrazu, a $z_0$ jest współrzędną piksela.
Gdy $|z_n|^2 > 4$, ciąg ucieka do nieskończoności i punkt nie należy do zbioru.
Porównujemy kwadrat modułu, więc nie liczymy pierwiastka.

\textbf{Piksel i płaszczyzna zespolona.} Piksel $(p_x, p_y)$ odpowiada
punktowi $x = c_x + (p_x - W/2)\,s$, $y = c_y - (p_y - H/2)\,s$, gdzie $s$ to
skala. Przy przybliżaniu o czynnik $k$ nowy środek wynosi
$\mathbf{m}' = \mathbf{q} + (\mathbf{m} - \mathbf{q})\,k$, gdzie $\mathbf{q}$
to punkt pod kursorem. Dzięki temu punkt pod kursorem zostaje w miejscu.

\textbf{Płynne kolorowanie.} Zamiast całkowitej liczby kroków $n$ używamy
wartości ułamkowej
\[
  \nu = n + 1 - \log_2\bigl(\ln |z_n|\bigr),
\]
którą przeliczamy na kolor z palety. Usuwa to pasy kolorów. Punkty należące
do zbioru są czarne.

\subsection{Fraktale IFS -- gra w chaos}
\label{sec:ifs}

Zaczynamy od dowolnego punktu. W każdym kroku losujemy przekształcenie $w_i$
z prawdopodobieństwem $p_i$, stosujemy je do punktu i zaznaczamy wynik.
Pierwsze 20 punktów pomijamy, bo nie leżą jeszcze na atraktorze.

Prawdopodobieństwa $p_i$ pochodzą z pliku albo z panelu. Przy ręcznej edycji
panel podpowiada wartości proporcjonalne do pola, jakie przekształcenie
zachowuje:
\[
  p_i = \frac{\max\bigl(|a_i d_i - b_i c_i|,\ \varepsilon\bigr)}
             {\sum_j \max\bigl(|a_j d_j - b_j c_j|,\ \varepsilon\bigr)} .
\]
Przy takich wartościach punkty rozkładają się na fraktalu równomiernie.
Minimalna wartość $\varepsilon$ chroni przekształcenia spłaszczające,
np.\ łodygę paproci.

\textbf{Dopasowanie do okna.} Najpierw liczymy kilka tysięcy punktów, aby
wyznaczyć prostokąt otaczający atraktor. Dodajemy margines 5\% i dobieramy
skalę z zachowaniem proporcji, tak aby fraktal wypełniał okno.

\textbf{Kolorowanie.} Każde przekształcenie ma własny kolor, a punkt
przyjmuje kolor przekształcenia, które go wyznaczyło. Opcjonalnie zliczamy
trafienia w piksel i ustalamy jasność logarytmicznie:
$\mathrm{jasność} = \ln(1 + h) / \ln(1 + h_{\max})$.

\subsection{Animacja przejścia między fraktalami IFS}

Niewielka zmiana współczynników przekształceń daje niewielką zmianę obrazu.
Dlatego klatkę $k = 0, 1, \ldots, K-1$ przejścia od fraktala $A$ do fraktala
$B$ tworzymy, zmieniając liniowo każdy współczynnik:
\[
  w_i^{(k)} = (1 - t)\, w_i^{A} + t\, w_i^{B},
  \qquad t = \frac{k}{K-1} \in [0, 1] .
\]
Pierwsza klatka przedstawia więc fraktal $A$, a ostatnia fraktal $B$.

\begin{itemize}
  \item \textbf{Różna liczba przekształceń.} Do zbioru z mniejszą liczbą
        przekształceń dokładamy kopie istniejących, dzieląc prawdopodobieństwo
        kopiowanego przekształcenia na pół. Dwie identyczne kopie dają ten sam
        atraktor, więc pierwsza klatka się nie zmienia.
  \item \textbf{Poprawność klatek pośrednich.} Jeśli przekształcenia $A$ i $B$
        są zwężające, to ich kombinacja wypukła też jest zwężająca:
        $\lVert (1-s)M_A + s M_B \rVert \le (1-s)\lVert M_A \rVert
        + s\lVert M_B \rVert < 1$. Każda klatka jest więc poprawnym fraktalem.
  \item \textbf{Stabilny kadr.} Prostokąt otaczający liczymy dla każdej klatki
        i wygładzamy go średnią kroczącą, aby obraz nie ,,pływał''.
  \item \textbf{Eksport.} Każde wywołanie \kod{tick()} liczy jedną klatkę,
        zapisuje ją do pliku \kod{klatka\_NNNN.png} (lub \kod{.bmp})
        i pokazuje w oknie. Między klatkami interfejs obsługuje zdarzenia,
        więc działa podgląd, pasek postępu i przycisk \textit{Przerwij}.
        W poszerzeniu klatki liczy wątek roboczy, a przerwanie sygnalizuje
        flaga \kod{std::atomic<bool>}.
\end{itemize}

\subsection{L-systemy i grafika żółwia}

L-system składa się z aksjomatu, reguł przepisywania i kąta. W każdej
iteracji wszystkie symbole napisu zastępujemy jednocześnie według reguł.
Długość napisu rośnie wykładniczo, więc ograniczamy ją, np.\ do 5~mln
symboli, i ostrzegamy użytkownika.

\begin{center}
\begin{tabular}{@{}cl@{}}
  \toprule
  Symbol & Polecenie dla żółwia \\
  \midrule
  \kod{F} & idź naprzód o krok i rysuj odcinek \\
  \kod{f} & idź naprzód o krok bez rysowania \\
  \kod{+}, \kod{-} & obróć się w lewo lub w prawo o zadany kąt \\
  \kod{[} & zapamiętaj pozycję i kierunek (odłóż na stos) \\
  \kod{]} & wróć do ostatnio zapamiętanego stanu (zdejmij ze stosu) \\
  \bottomrule
\end{tabular}
\end{center}

Żółw przechodzi napis dwa razy. Za pierwszym razem tylko wyznacza prostokąt
otaczający rysunek, aby dobrać skalę. Za drugim rysuje odcinki algorytmem
Bresenhama bezpośrednio w buforze pikseli.

\subsection{Fraktale IFS 3D (poszerzenie)}

Przekształcenie przestrzenne ma 12 współczynników:
$\mathbf{x}' = M\mathbf{x} + \mathbf{t}$, gdzie $M$ to macierz $3 \times 3$.
Grę w chaos prowadzimy w przestrzeni. Każdy punkt przekształcamy do układu
kamery i rzutujemy perspektywicznie: $x_s = f\,x_c / z_c + W/2$. Własny bufor
głębi przechowuje dla piksela odległość najbliższego punktu, więc dalsze
punkty nie przesłaniają bliższych. Jasność maleje z odległością od
obserwatora.

\subsection{Obliczenia wielowątkowe (poszerzenie)}

Obraz dzielimy na poziome pasy, po jednym na wątek
(\kod{std::thread::hardware\_concurrency()}). Każdy wątek zapisuje tylko swoje
wiersze, więc nie potrzeba blokad. W grze w chaos każdy wątek ma własny
generator liczb losowych i własny bufor trafień, które na końcu sumujemy.
Obraz liczony na wielu wątkach musi być identyczny z obrazem liczonym na
jednym, co sprawdzamy w testach.

% ---------------------------------------------------------------
\section{Podział pracy i analiza czasowa}

Dokumentację i testy prowadzimy wspólnie. Na początku wspólnie ustalamy
część wspólną (\kod{Image}, \kod{LineReader}, \kod{ParseError},
\kod{FractalRenderer}) i szkielet aplikacji z zaślepkami trzech trybów.
Od tej chwili każdy pracuje w swoim module i widzi wynik w programie.
Tygodnie liczymy od akceptacji tematu.

\begin{center}
\renewcommand{\arraystretch}{1.25}
\begin{tabularx}{\textwidth}{@{}>{\raggedright\arraybackslash}p{2.2cm}
  >{\raggedright\arraybackslash}X >{\raggedright\arraybackslash}X
  >{\raggedright\arraybackslash}X@{}}
  \toprule
  Etap & Jakub Piela & Marcel Gładysz & Patryk Pilch \\
  & \footnotesize interfejs, zbiory zespolone & \footnotesize IFS i animacja
  & \footnotesize L-systemy \\
  \midrule
  \textbf{Tydz. 1--2}\newline szkielet &
    szkielet aplikacji (CMake, okno, panel, menu, przełączanie trybów), nawigacja myszą, odwzorowanie piksel--płaszczyzna, Mandelbrot &
    parser \kod{.ifs}, gra w chaos &
    parser \kod{.lsys}, przepisywanie napisów \\
  \textbf{Tydz. 3--4}\newline rdzeń &
    Julia (wybór $c$ kliknięciem), płynne kolorowanie, palety &
    wbudowane fraktale (Sierpiński, paproć), edycja współczynników, autoskalowanie, kolorowanie &
    grafika żółwia z~rozgałęzieniami, wbudowane krzywe, dopasowanie do okna \\
  \textbf{Tydz. 5--6}\newline komplet podstaw &
    zapis obrazu PNG/BMP, okna dialogowe, integracja modułów &
    animacja przejścia dwóch fraktali, podgląd klatek, eksport ponumerowanych klatek &
    wczytywanie i~zapis definicji, testy parserów \\
  \textbf{Tydz. 7--8}\newline poszerzenia &
    wielowątkowość i~animacja w~tle z~paskiem postępu, animacja przybliżania, zapis w~wysokiej rozdzielczości, archiwum 7z &
    plik scenariusza, fraktale IFS 3D &
    kolejne zbiory zespolone (Burning Ship, Multibrot, Newton), edytor palet \\
  \bottomrule
\end{tabularx}
\end{center}

Wymagania podstawowe mają być gotowe po szóstym tygodniu. Ostatnie dwa
tygodnie zostają na poszerzenia, testy i dokumentację w PDF.

% ---------------------------------------------------------------
\section{Plan testów}

Testy modułów działają bez okna i uruchamiamy je poleceniem \kod{ctest}.

\begin{itemize}
  \item \textbf{Testy modułów:} parsery na poprawnych i celowo błędnych
        plikach (oczekiwany komunikat z numerem wiersza); znane wyniki: punkt
        $c = 0$ należy do zbioru Mandelbrota, a $c = 1$ ucieka po 3 krokach;
        liczba symboli \kod{F} w krzywej Kocha po $n$ iteracjach wynosi $5^n$;
        współczynniki pierwszej i ostatniej klatki przejścia są równe
        współczynnikom fraktali $A$ i $B$.
  \item \textbf{Testy powiązań:} zmiana trybu w trakcie generowania animacji,
        zapis obrazu po wczytaniu definicji, przerwanie eksportu animacji.
  \item \textbf{Testy całościowe:} przejście trójkąta Sierpińskiego w paproć
        Barnsleya w 100 klatkach, zmiana rozmiaru okna, pomiar czasu obliczeń
        klatki.
  \item \textbf{Niezmienniki:} liczba zapisanych klatek równa zadanej liczbie
        klatek przejścia; w poszerzeniu obraz zbioru zespolonego nie zależy od
        liczby wątków.
\end{itemize}

% ---------------------------------------------------------------
\section{Ryzyka i plan awaryjny}

\begin{itemize}
  \item \textbf{Wydajność przy dużym przybliżeniu} -- w wersji podstawowej
        liczymy w jednym wątku, więc podczas przeciągania pokazujemy podgląd
        w niższej rozdzielczości, a pełny obraz po puszczeniu przycisku myszy.
        Obliczenia wielowątkowe są pierwszym poszerzeniem.
  \item \textbf{Precyzja \kod{double}} -- przybliżanie ograniczamy do zakresu,
        w którym \kod{double} wystarcza, i nie pozwalamy przybliżać dalej.
  \item \textbf{Różne systemy w zespole} -- budowanie przez CMake, biblioteki
        jako submoduły w ustalonych wersjach, instrukcja budowania w
        \kod{README.md}.
  \item \textbf{Bardzo długie napisy L-systemów} -- limit długości
        i ostrzeżenie zamiast zawieszenia programu.
  \item \textbf{Zakres} -- jeśli zabraknie czasu, wymagania podstawowe
        pozostają nietknięte, a odpadają pozycje z listy poszerzeń, od końca.
\end{itemize}

\end{document}